\documentclass{article}
\usepackage{amsmath, amssymb, amsthm}
\usepackage{geometry}
\geometry{a4paper, margin=1in}

\title{Optimization HW1 Solutions}
\author{}
\date{}

\begin{document}

\maketitle

\section*{Problem 1: Convex Functions [20 pts]}

\subsection*{(a) $f(x) = \frac{1}{x^2}$ over domain $(0, \infty)$.}

\textbf{Answer: Strictly Convex}

\begin{proof}
To determine the convexity, we examine the second derivative of the function $f(x) = x^{-2}$.
\begin{align*}
f'(x) &= -2x^{-3} \\
f''(x) &= 6x^{-4} = \frac{6}{x^4}
\end{align*}
Since the domain is $(0, \infty)$, we have $x > 0$, which implies $x^4 > 0$.
Therefore, $f''(x) = \frac{6}{x^4} > 0$ for all $x \in (0, \infty)$.

Since the second derivative is strictly positive everywhere on the domain, the function is \textbf{strictly convex}.

It is not $\mu$-strongly convex because $\lim_{x \to \infty} f''(x) = 0$. There is no constant $\mu > 0$ such that $f''(x) \ge \mu$ for all $x$.
\end{proof}

\subsection*{(b) $f(x) = \min\{x, 1-x\}$ over domain $\mathbb{R}$.}

\textbf{Answer: Concave}

\begin{proof}
This function is the minimum of two affine (linear) functions: $g(x) = x$ and $h(x) = 1-x$.

A known property of convex/concave functions is that the pointwise minimum of a set of concave functions is itself concave. Since affine functions are both convex and concave, $g(x)$ and $h(x)$ are concave. Therefore, their minimum $f(x)$ is \textbf{concave}.

\textbf{Geometric Intepretation:}
The function looks like an inverted "V" with the peak at $x = 0.5$ (where $x = 1-x$).
\begin{itemize}
    \item For $x \le 0.5$, $f(x) = x$ (slope 1).
    \item For $x > 0.5$, $f(x) = 1-x$ (slope -1).
\end{itemize}
It is not strictly concave because it is linear (affine segments) on the intervals $x < 0.5$ and $x > 0.5$. It is not strongly concave for the same reason (second derivative is 0 almost everywhere).
\end{proof}

\subsection*{(c) A McCulloch-Pitts neuron with ReLU activation: $f(x) = (\theta^T x)_+$ over domain $\mathbb{R}^n$.}

\textbf{Answer: Convex}

\begin{proof}
The function is given by $f(x) = \max\{0, \theta^T x\}$.
This is a composition of the function $g(y) = \max\{0, y\}$ (ReLU) and the affine mapping $h(x) = \theta^T x$.

\begin{enumerate}
    \item \textbf{Convexity of ReLU:} The function $g(y) = \max\{0, y\}$ is convex. It is the maximum of two convex functions ($y$ and $0$).
    \item \textbf{Composition Rule:} If $g: \mathbb{R} \to \mathbb{R}$ is convex and non-decreasing, and $h: \mathbb{R}^n \to \mathbb{R}$ is convex, then $g(h(x))$ is convex. Alternatively, a more general rule applies simply because the inner function is affine: If $g$ is convex and $h$ is affine, then $g(h(x))$ is convex.
\end{enumerate}

Since $g(y)$ is convex and $x \mapsto \theta^T x$ is affine, the composition $f(x)$ is \textbf{convex}.

It is not strictly convex because on the half-space where $\theta^T x \le 0$, the function is constant ($f(x) = 0$), i.e., "flat".
\end{proof}

\subsection*{(d) A McCulloch-Pitts neuron with sigmoid activation: $f(x) = \text{sigm}(\theta^T x)$ over domain $\mathbb{R}^n$.}

\textbf{Answer: Neither convex nor concave}

\begin{proof}
Let $g(z) = \frac{1}{1+e^{-z}}$ be the sigmoid function. Then $f(x) = g(\theta^T x)$.
Since $\theta^T x$ is an affine mapping, the convexity of $f(x)$ depends entirely on the convexity of the scalar sigmoid function $g(z)$ along the line defined by $\theta$.

Let's examine $g''(z)$:
\begin{align*}
g'(z) &= g(z)(1 - g(z)) \\
g''(z) &= g'(z)(1 - g(z)) + g(z)(-g'(z)) \\
&= g(z)(1 - g(z))(1 - 2g(z))
\end{align*}

The sign of $g''(z)$ depends on the sign of $(1 - 2g(z))$.
Since $g(z) \in (0, 1)$:
\begin{itemize}
    \item If $g(z) < 1/2$ (which corresponds to $z < 0$), then $1 - 2g(z) > 0 \implies g''(z) > 0$. The function is convex in this region.
    \item If $g(z) > 1/2$ (which corresponds to $z > 0$), then $1 - 2g(z) < 0 \implies g''(z) < 0$. The function is concave in this region.
\end{itemize}

Since the domain is $\mathbb{R}^n$ (and assuming $\theta \neq 0$), the argument $\theta^T x$ takes values in $(-\infty, \infty)$. The function transitions from convex to concave. Therefore, it is \textbf{neither convex nor concave} globally.
\end{proof}

\end{document}
